(1) 两个恒温热浴温度 \(T_h>T_c>0\)，热机内部吸放热端温度为 \(\vartheta_h,\vartheta_c\)，接触热导分别为 \(K_h,K_c>0\)，稳态热流为 \(Q_h=K_h(T_h-\vartheta_h)\)、\(Q_c=K_c(\vartheta_c-T_c)\)。若内部循环可逆，建立熵流约束。

(2) 以内部共同熵流率 \(s=Q_h/\vartheta_h=Q_c/\vartheta_c\) 为变量，求两端工作温度、功率 \(P(s)\) 及允许范围。

(3) 对 \(s\) 优化功率，求最大功率熵流和最大功率效率，并核验其低于Carnot效率。

(4) 求两外部热浴总熵产生率，证明在非零输出功率时严格为正；讨论接触热导无限大时的可逆极限。
